\section{The record: from case files to dated events}\label{sec:record}

\subsection{Corpus and parse}

The USPTO Trademark Annual Applications backfile (product TRTYRAP) is the
complete record of US trademark filings from 1884 onward: 13.99 million
case files, each carrying the goods/services description by class, the
prosecution history, and the ownership and correspondence records. The
prosecution history is a list of coded events with dates. This paper
re-parses every case file into 242 million dated events, keyed by serial
number, so that for each application the sequence of office actions,
responses, allowances, statements of use, registration, Section~8 and
Section~9 filings, cancellations and renewals is available with its
date.\footnote{Event codes are read against the USPTO's own event-code
dictionary. Terminal status codes are retained for cross-checking: 700
registered; 701--705 Section~8 accepted; 710 cancelled for Section~8
failure; 800 renewed; 900 expired. The year-six cancellation and the
year-ten cancellation share status 710 and are separable only by event
date, which is the reason the corpus is event-dated.} From the same files
the parse reads counsel of record on each filing (and, for refilings, on
each attempt separately), the statutory filing basis (use, intent to use,
\S44(e) home registration, \S66(a) Madrid extension), the owner's
domicile, and the owner name, which is normalized --- uppercased,
stripped of punctuation and of two dozen legal suffixes --- so that
\textsc{acme corp.} and \textsc{acme corporation} are one owner.
\citet{Graham2013} describe the administrative data; the event-dated
parse adds the full prosecution timeline of each case.

Response latency in prosecution is a by-product of the parse and is
retained as a health indicator: the median days from an office action to
the applicant's response, per filing.
